% ---
% id: meat-core/frank-wolfe/req:returned-point-measured
% kind: req
% title: "Measurement of the returned point"
% summary: >-
%   Every run applies the three iterate checks once more to the point it
%   returned and records whether that point passes, with the three quantities
%   measured. Neither the bound, nor the tier, nor the key length rests on
%   that measurement.
% prominence: standard
% status: implemented
% uses: [thm:firstorder, req:iterate-checks]
% mentions: [rem:fo-returned-point]
% assumes: []
% refs: []
% provenance: newly-derived
% implements: certiqs_key.solvers.frank_wolfe.joint_driver:run_joint_frank_wolfe
% lean:
% ---

The checks of Requirement~\ref{req:iterate-checks} are made on the point
consumed at each iterate, and the point a run returns, at which its
\tref{estimate} is taken, need not be one of those: a run that stops at its
iteration limit returns a point formed after the last point it consumed.
What the measurement establishes of that point, the $J$ factor of the
estimate's premise and no more, is Remark~\ref{rem:fo-returned-point}.

\begin{requirement}\label{req:returned-point-measured}
Every run applies the three checks of Requirement~\ref{req:iterate-checks}
once more, at the end of the run, to the point it returned: per-block
positivity at exactly $0$, the coupling residual at \iface{iterate_coupling}
and each $\bowtie_k$ residual at \iface{constraint_rows}. It records whether
that point passes all three, together with the three quantities measured,
in \iface{optimizer_certificate}. The measurement enters neither the
first-order bound, nor the tier, nor the key length.
\end{requirement}
